% Part: first-order-logic
% Chapter: syntax-and-semantics
% Section: substitution

\documentclass[../../../include/open-logic-section]{subfiles}

\begin{document}

\olfileid{fol}{syn}{ext}

\olsection{విస్తారత}

\begin{explain}
విస్తారతను కొన్నిసార్లు ప్రాసంగికత అని కూడా అంటారు. దీన్ని అనౌపచారికంగా ఇలా
చెప్పవచ్చు: \tetoken{ఒక నిర్మాణం}{!!a{structure}}~$\Struct M$లో, ఒక \tetoken{చర}{!!a{variable}} నిర్దేశం~$s$కు
సాపేక్షంగా \tetoken{సూత్రం}{!!{formula}}~$!A$ సంతృప్తమవుతుందో లేదో నిర్ణయించేవి
\tetoken{వ్యక్తి క్షేత్రపు}{!!{domain}} పరిమాణం, అలాగే $!A$లో కనిపించే భాషా సంకేతాలకు
$\Struct M$, ~$s$లు ఇచ్చే నిర్దేశాలు మాత్రమే.

విస్తారత నుంచి నేరుగా ఒక పర్యవసానం వస్తుంది. ఒక వాక్యం~$!A$లో కనిపించే భాషా
సంకేతాలన్నింటికీ రెండు \tetoken{నిర్మాణాలు}{!!{structure}s}~$\Struct M$, ~$\Struct M'$లు
ఒకే అర్థాన్ని ఇచ్చి, వాటి వ్యక్తి క్షేత్రం ఒకటే అయితే, $!A$ సత్యమా కాదా అనే
విషయంలో కూడా $\Struct M$, ~$\Struct M'$లు ఒకే ఫలితాన్ని ఇస్తాయి.
\end{explain}

\begin{prop}[విస్తారత]
  \ollabel{prop:extensionality}
  $!A$ ఒక \tetoken{సూత్రం}{!!a{formula}} అనీ, $\Struct M_1$, ~$\Struct M_2$లు
  $\Domain{M_1} = \Domain{M_2}$ గల \tetoken{నిర్మాణాలు}{!!{structure}s} అనీ, ~$s$ అనేది
  $\Domain{M_1} = \Domain{M_2}$పై నిర్వచించిన చర నిర్దేశమనీ అనుకుందాం. $!A$లో
  కనిపించే ప్రతి \tetoken{స్థిరసంకేతం}{!!{constant}}~$c$, సంబంధ సంకేతం~$R$, \tetoken{ప్రమేయ సంకేతం}{!!{function}}~$f$కు
  $\Assign{c}{M_1} = \Assign{c}{M_2}$,
  $\Assign{R}{M_1}=\Assign{R}{M_2}$, అలాగే
  $\Assign{f}{M_1} = \Assign{f}{M_2}$ అయితే,
  $\Sat{M_1}{!A}[s]$ అవుతుంది అప్పుడే, అప్పుడు మాత్రమే
  $\Sat{M_2}{!A}[s]$.
\end{prop}

\begin{proof}
  మొదట, ఈ సూత్రంలో కనిపించే తార్కికేతర సంకేతాలనే వాడే ప్రతి పదానికీ, ఏ ఉమ్మడి చర నిర్దేశానికైనా
  $\Value{t}{M_1}[s] = \Value{t}{M_2}[s]$ అని ($t$పై ఆగమనం ద్వారా) నిరూపించండి. తరువాత అదే సంకేతాలున్న భాషలో
  $!A$ పరమాణువుగా ఉండే ఆధార సందర్భంలో పదాలకు నిరూపించిన ఫలితాన్ని వాడి, ~$!A$పై ఆగమనం ద్వారా
  ప్రతిపాదనను నిరూపించండి.
\end{proof}

\begin{prob}
\olref[fol][syn][ext]{prop:extensionality} నిరూపణను వివరంగా
రాయండి.
\end{prob}

\begin{cor}[\tetoken{వాక్యాలకు}{!!^{sentence}s} విస్తారత]
  \ollabel{cor:extensionality-sent}
  $!A$ ఒక \tetoken{వాక్యం}{!!a{sentence}} అనీ, $\Struct{M_1}$, ~$\Struct{M_2}$లు
  \olref{prop:extensionality}లో ఉన్నట్లే అనీ అనుకుందాం. అప్పుడు
  $\Sat{M_1}{!A}$ అవుతుంది అప్పుడే, అప్పుడు మాత్రమే $\Sat{M_2}{!A}$.
\end{cor}

\begin{proof}
\olref[ass]{cor:sat-sentence}ను ఉపయోగిస్తే ఈ ఫలితం
\olref{prop:extensionality} నుంచి వస్తుంది.
\end{proof}

అంతేకాక, ఒక పదానికి గల విలువా, \tetoken{ఒక నిర్మాణం}{!!a{structure}} ఒక \tetoken{సూత్రాన్ని}{!!a{formula}}
సంతృప్తిపరుస్తుందా లేదా అన్నదీ, వాటి ఉపపదాల విలువలపైన మాత్రమే
ఆధారపడతాయి.

\begin{prop}\ollabel{prop:ext-terms}
$\Struct M$ ఒక \tetoken{నిర్మాణం}{!!a{structure}}, ~$t$, ~$t'$లు పదాలు, ~$s$ ఒక చర
నిర్దేశం అనుకుందాం. అప్పుడు $\Value{\Subst{t}{t'}{x}}{M}[s] =
\Value{t}{M}[\Subst{s}{\Value{t'}{M}[s]}{x}]$.
\end{prop}

\begin{proof}
$t$పై ఆగమనం ద్వారా నిరూపిస్తాం.
\begin{enumerate}
\item $t$ ఒక స్థిరసంకేతం, అంటే $t\ident c$ అయితే,
  $\Subst{t}{t'}{x} = c$; ఇంకా $\Value{c}{M}[s] = \Assign{c}{M} =
  \Value{c}{M}[\Subst{s}{\Value{t'}{M}[s]}{x}]$.

\item $t$ అనేది $x$ కాని చరం, అంటే $t \ident y$ అయితే,
  $\Subst{t}{t'}{x} = y$; ఇంకా
  $\varAssign{s}{\Subst{s}{\Value{t'}{M}[s]}{x}}{x}$ కాబట్టి
  $\Value{y}{M}[s] =
  \Value{y}{M}[\Subst{s}{\Value{t'}{M}[s]}{x}]$.

\item $t \ident x$ అయితే, $\Subst{t}{t'}{x} = t'$. కానీ
  $\Subst{s}{\Value{t'}{M}[s]}{x}$ నిర్వచనం ప్రకారం
  $\Value{x}{M}[\Subst{s}{\Value{t'}{M}[s]}{x}] = \Value{t'}{M}[s]$.

\item $t \ident \Atom{f}{t_1,\dots,t_n}$ అయితే మనకు:
\begin{multline*}
  \Value{\Subst{t}{t'}{x}}{M}[s] \\
  \begin{aligned}[b]
& = \Value{\Atom{f}{\Subst{t_1}{t'}{x}, \dots, \Subst{t_n}{t'}{x}}}{M}[s]\\
& \qquad    \text{$\Subst{t}{t'}{x}$ నిర్వచనం ప్రకారం}\\
& = \Assign{f}{M}(\Value{\Subst{t_1}{t'}{x}}{M}[s], \dots,
    \Value{\Subst{t_n}{t'}{x}}{M}[s])\\
    & \qquad  \text{$\Value{\Atom{f}{\dots}}{M}[s]$ నిర్వచనం ప్రకారం}\\
& = \Assign{f}{M}(\Value{t_1}{M}[\Subst{s}{\Value{t'}{M}[s]}{x}], \dots,
   \Value{t_n}{M}[\Subst{s}{\Value{t'}{M}[s]}{x}])\\
& \qquad    \text{ఆగమన పరికల్పన ప్రకారం}\\
& = \Value{t}{M}[\Subst{s}{\Value{t'}{M}[s]}{x}]
    \text{$\Value{\Atom{f}{\dots}}{M}[\Subst{s}{\Value{t'}{M}[s]}{x}]$ నిర్వచనం ప్రకారం}
  \end{aligned}
\end{multline*}
\end{enumerate}
\end{proof}

\begin{prop}\ollabel{prop:ext-formulas} $\Struct M$ ఒక
\tetoken{నిర్మాణం}{!!a{structure}}, ~$!A$ ఒక \tetoken{సూత్రం}{!!a{formula}}, ~$t'$ అనేది $!A$లో $x$కు బదులుగా ప్రతిస్థాపించడానికి స్వేచ్ఛగా ఉండే పదం, ~$s$ ఒక చర
నిర్దేశం అనుకుందాం. అప్పుడు $\Sat{M}{\Subst{!A}{t'}{x}}[s]$ అవుతుంది
అప్పుడే, అప్పుడు మాత్రమే
$\Sat{M}{!A}[\Subst{s}{\Value{t'}{M}[s]}{x}]$.
\end{prop}

\begin{proof}
సాధన.
\end{proof}

\begin{prob}
\olref[fol][syn][ext]{prop:ext-formulas}ను నిరూపించండి.
\end{prob}

\begin{explain}
  \cref{fol:syn:ext:prop:ext-terms,fol:syn:ext:prop:ext-formulas}కు సంబంధించిన
  విషయం ఇది. మనకు ఒక పదం~$t$ లేదా \tetoken{ఒక సూత్రం}{!!a{formula}}~$!A$, మరొక పదం~$t'$
  ఉన్నాయనీ, $\Subst{t}{t'}{x}$ విలువను లేదా
  $\Subst{!A}{t'}{x}$ అనేది \tetoken{ఒక నిర్మాణం}{!!a{structure}}~$\Struct M$లో
  ఒక \tetoken{చర}{!!a{variable}} నిర్దేశం~$s$కు సాపేక్షంగా సంతృప్తమవుతుందో లేదో
  తెలుసుకోవాలనీ అనుకుందాం. అప్పుడు ముందుగా ప్రతిస్థాపన చేసి, తరువాత
  $\Struct{M}$, ~$s$లకు సాపేక్షంగా విలువను లేదా సంతృప్తిని పరిశీలించవచ్చు.
  లేదా ముందుగా $\Struct{M}$లో $s$కు సాపేక్షంగా $t'$ విలువ
  $m = \Value{t'}{M}[s]$ను నిర్ణయించి, \tetoken{చర}{!!{variable}} నిర్దేశాన్ని
  $\Subst{s}{m}{x}$గా మార్చి, తరువాత $\Struct{M}$లో
  $\Subst{s}{m}{x}$కు సాపేక్షంగా~$t$ విలువను, లేదా
  $\Sat{M}{!A}[\Subst{s}{m}{x}]$ అవుతుందో పరిశీలించవచ్చు.
  \Cref{fol:syn:ext:prop:ext-terms,fol:syn:ext:prop:ext-formulas} ప్రకారం,
  ఏ పద్ధతి వాడినా సమాధానం ఒకటే.
\end{explain}

\end{document}
